% !TEX root = ../main.tex

\section{Sistemi di Riferimento}
Per descrivere la dinamica del tricottero VTOL, si adottano due sistemi di riferimento principali, entrambi destrorsi:
\begin{itemize}
    \item \textbf{Sistema Inerziale (Earth Frame) $\mathcal{F}_g$}: Fisso rispetto al terreno, definito dalla terna $(O_g, X_g, Y_g, Z_g)$. La posizione del velivolo è indicata dal vettore $P_g = [x_g, y_g, z_g]^T$ e l'orientamento dagli angoli di Eulero $\Theta_g = [\phi, \theta, \psi]^T$ (Roll, Pitch, Yaw).
    \item \textbf{Sistema Solidale (Body Frame) $\mathcal{F}_b$}: Con origine nel centro di massa (CoM) del velivolo e assi $(X_b, Y_b, Z_b)$ solidali alla struttura. Le velocità lineari e angolari sono definite rispettivamente come $V_b = [u, v, w]^T$ e $\Omega_b = [p, q, r]^T$.
\end{itemize}
\begin{figure}[h]
    \centering
    \includegraphics[width=0.7\textwidth]{sist_rif.png}
    \caption{Sistemi di riferimento.}
    \label{fig:sist_rif}
\end{figure}
La trasformazione delle velocità lineari dal sistema solidale al corpo a quello inerziale avviene tramite la matrice di rotazione $R_{b \to g}$ (convenzione ZYX):
\begin{equation}
    \dot{P}_g = R_{b \to g}(\Theta_g) V_b
\end{equation}
Analogamente, la relazione cinematica tra le velocità angolari del corpo e le derivate degli angoli di Eulero è data dalla matrice di trasformazione $J_{b \to g}$:
\begin{equation}
    \dot{\Theta}_g = J_{b \to g}(\Theta_g) \Omega_b
\end{equation}

\section{Equazioni della Dinamica}
Il modello dinamico è derivato utilizzando la formulazione di Newton-Eulero per un corpo rigido a 6 gradi di libertà. Le equazioni vettoriali che descrivono il moto traslazionale e rotazionale sono:

\begin{equation}
    \begin{cases}
        m \dot{V}_b + \Omega_b \times (m V_b) = F_{\mathrm{tot}} \\
        I_b \dot{\Omega}_b + \Omega_b \times (I_b \Omega_b) = M_{\mathrm{tot}}
    \end{cases}
    \label{eq:newton_euler}
\end{equation}
\begin{itemize}
    \item $m$ è la massa totale del velivolo,
    \item $I_b$ è il tensore d'inerzia (assunto diagonale per simmetria strutturale), 
    \item $F_{\mathrm{tot}}$ e $M_{\mathrm{tot}}$ sono rispettivamente la somma delle forze e dei momenti esterni agenti sul corpo, espressi nel \textit{body frame}.

\end{itemize} 

\section{Analisi delle Forze}
Le forze agenti sul tricottero sono la risultante della spinta dei motori ($F_{\mathrm{th}}$), della forza di gravità ($F_g$) e delle forze aerodinamiche ($F_{\mathrm{aero}}$):
\begin{equation}
    F_{\mathrm{tot}} = F_{\mathrm{th}} + F_g + F_{\mathrm{aero}}
\end{equation}

\subsection{Thrust}
Il velivolo è dotato di tre rotori: due anteriori (1 e 2) e uno posteriore (3).
\begin{itemize}
    \item I rotori anteriori possono inclinarsi (\textit{tilt}) attorno all'asse $Y_b$ di un angolo $\theta_1$ e $\theta_2$.
    \item Il rotore di coda possiede due gradi di libertà, potendo ruotare attorno a $Y_b$ ($\theta_3$) e $Z_b$ ($\theta_4$).
\end{itemize}
Ogni rotore genera una spinta $f_i = k \omega_i^2$ lungo il proprio asse di rotazione. Il vettore forza totale, espresso nel body frame, si ottiene sommando i contributi dei tre rotori ruotati dalle rispettive matrici di tilt:
\begin{figure}[h]
    \centering
    \includegraphics[width=0.6\textwidth]{rotori_1_2.png}
    \caption{Direzione (in rosso) del thrust per i rotori 1 e 2 nel piano $(X_{b}, Z_{b})$.}
    \label{fig:rotori_1_2}
\end{figure}

\begin{figure}[H]
    \centering
    \includegraphics[width=0.7\textwidth]{rotore_3.png}
    \caption{Direzione (in rosso) del thrust per il rotore 3 nei piani $(X_{b}, Z_{b})$ e $(X_{b}, Y_{b})$.}
    \label{fig:rotore_3}
\end{figure}

\begin{equation}
    F_{\mathrm{th}} = \sum_{i=1}^{3} R_{\mathrm{tilt}, i} \cdot \begin{bmatrix} 0 \\ 0 \\ -f_i \end{bmatrix}
\end{equation}
Espandendo i termini, la forza risultante assume la forma:
\begin{equation}
    F_{\mathrm{th}} = 
    \begin{bmatrix}
    c_{\theta_1} f_1 + c_{\theta_2} f_2 + c_{\theta_3} c_{\theta_4} f_3 \\
    c_{\theta_3} s_{\theta_4} f_3 \\
    -s_{\theta_1} f_1 - s_{\theta_2} f_2 - s_{\theta_3} f_3
    \end{bmatrix}
\end{equation}
dove si è usata la notazione abbreviata $c_\alpha = \cos\alpha$ e $s_\alpha = \sin\alpha$.

\subsection{Gravità e Aerodinamica}
La forza peso, costante nel sistema inerziale, viene ruotata nel sistema solidale al corpo:
\begin{equation}
    F_g = R_{b \to g}^T \begin{bmatrix} 0 \\ 0 \\ mg \end{bmatrix} = mg \begin{bmatrix} -\sin\theta \\ \cos\theta \sin\phi \\ \cos\theta \cos\phi \end{bmatrix}
\end{equation}
Le forze aerodinamiche (Portanza $L$ e Resistenza $D$) agenti sulle superfici alari sono modellate in funzione della pressione dinamica e dei coefficienti aerodinamici $C_L$ e $C_D$:
\begin{equation}
    F_{\mathrm{aero}} \approx 
    \begin{bmatrix}
        -\frac{1}{2}\rho S C_D v_x |v_x| \\
        0 \\
        -\frac{1}{2}\rho S C_L v_x^2 - \frac{1}{2}\rho S C_{D_z} v_z |v_z|
    \end{bmatrix}
\end{equation}
Si assume per semplicità che il moto laterale non generi forze aerodinamiche significative sulle ali principali.

\section{Analisi dei Momenti}
Il momento totale agente sul baricentro è la somma di diverse componenti:
\begin{equation}
    M_{\mathrm{tot}} = M_{\mathrm{th}} + M_{\mathrm{drag}} + M_{\mathrm{aero}} + M_{\mathrm{gyro}} + M_{\mathrm{fin}}
\end{equation}

\subsection{Momenti generati dai Rotori}
Il contributo principale è dato dal momento generato dalla spinta dei rotori, calcolato come il prodotto vettoriale tra i bracci di leva $r_i$ (distanza del rotore dal CoM) e le forze di spinta $F_{\mathrm{th},i}$:
\begin{equation}
    M_{\mathrm{th}} = \sum_{i=1}^{3} (r_i \times F_{\mathrm{th},i})
\end{equation}
Inoltre, la rotazione delle eliche genera un momento torcente ($M_{\mathrm{drag}}$) opposto al verso di rotazione, proporzionale al quadrato della velocità angolare ($\tau_i = b \omega_i^2$). Poiché i rotori sono inclinabili, anche questo momento viene proiettato sugli assi del body frame in funzione degli angoli di tilt.

\subsection{Effetti Aerodinamici e Giroscopici}
\begin{itemize}
    \item \textbf{Momento Aerodinamico ($M_{\mathrm{aero}}$):} È generato dalle forze di portanza e resistenza applicate nel centro di pressione delle ali, che non coincide con il centro di massa. Questo crea coppie di beccheggio e rollio in funzione della velocità di avanzamento.
    \item \textbf{Momento Giroscopico ($M_{\mathrm{gyro}}$):} Dovuto alla precessione giroscopica dei rotori ad alta velocità quando il velivolo (o i rotori stessi tramite il meccanismo di tilt) ruota. È dato da:
    \begin{equation}
        M_{\mathrm{gyro}} = \sum_{i=1}^{3} I_{\mathrm{rot}} (\Omega_b \times \hat{e}_{\mathrm{rot},i}) \omega_i
    \end{equation}
    \item \textbf{Momento della Pinna ($M_{\mathrm{fin}}$):} La superficie verticale di coda genera un momento stabilizzante di imbardata e un momento di rollio indotto, modellati in funzione delle velocità angolari $p, r$ e della velocità laterale.
\end{itemize}
L'insieme di queste equazioni definisce il modello matematico completo utilizzato per la sintesi del controllo nel Capitolo \ref{ch:controllo}.